% META: {"id": "TEXP-GRA-EX-019", "chapitre": "TEXP-GRAPHES", "type_objet": "exercice", "capacites": ["TEXP-GRA-C4"], "capacites_codes": ["C4"], "methodes": ["M2","M5"], "parcours": 2, "competences": ["modeliser","raisonner"], "duree_min": 25, "mode_creation": "ex_nihilo", "sources_inspiration": [], "fichier_tex": "chapitres/TEXP-GRAPHES/exercices/TEXP-GRA-EX-019.tex", "corrige_tex": "chapitres/TEXP-GRAPHES/corriges/TEXP-GRA-CO-019.tex", "status": "approved"}
% BEGIN-VERIFY
% from sympy import *
% x = symbols('x')
% C = x**3/3 - 5*x**2 + 30*x
% Cp = diff(C, x)
% assert expand(Cp) == x**2 - 10*x + 30
% # discriminant: 100 - 120 = -20 < 0, pas de racine reelle
% disc = 100 - 120
% assert disc < 0
% assert Cp.subs(x, 0) == 30
% assert Cp.subs(x, 5) == 5
% assert C.subs(x, 0) == 0
% END-VERIFY
\begin{exercice}{TEXP-GRA-EX-019}{3}{25}
Le cout de production de $x$ unites est modelise par $C(x) = \dfrac{x^3}{3} - 5x^2 + 30x$.
\begin{enumerate}
  \item Calculer le cout marginal $C'(x)$. Son signe permet-il d'identifier un cout minimum ?
  \item Calculer le cout moyen $\dfrac{C(x)}{x}$ et determiner s'il admet un minimum.
  \item Comparer les deux resultats.
\end{enumerate}

\end{exercice}
